\chapter{Nomenclatura}
In seguito vengono elencati, in ordine alfabetico, i termini principali utilizzati in questa relazione:
\begin{itemize}
    \item agente: attore che visualizza l'interfaccia grafica;
    \item cammino: insieme di nodi di un grafo che indica quali sono stati visitati e in che ordine;
    \item cammino hamiltoniano: cammino che passa per tutti i nodi del grafo una e una sola volta;
    \item cammino migliore: cammino di archi pesati la cui sommatoria dei pesi è minima;
    \item cammino minimo: vedi cammino migliore;
    \item cliente: cliente dell'azienda, rappresentato da un nodo di un grafo;
    \item cluster: insieme di elaboratori collegati in locale;
    \item grafo: struttura dati composta da un insieme di nodi e un insieme di archi, che collegano due nodi tra loro;
    \item grafo pesato: grafo in cui ogni arco ha un determinato peso, rappresentato da un valore;
    \item interfaccia grafica: modulo per la visualizzazione del risultato;
    \item nodo ACO: nodo del sistema che si occupa della computazione dell'algoritmo di Ant Colony Optimization;
    \item nodo utente: nodo del sistema che si occupa della gestione dei clienti;
    \item percorso: vedi cammino;
    \item peso di un arco: valore legato ad un arco del grafo;
    \item problema del commesso viaggiatore: problema in cui si vuole trovare il cammino hamiltoniano di costo minimo;
    \item sede: sede dell'azienda, dove può essere contenuto un cluster;
    \item ticket: ``biglietto'' virtuale per fare la gestione del grafo dei clienti in ordine;
    \item topologia del sistema: struttura dei collegamenti degli elementi del sistema;
    \item utente: colui che interagisce con le API del sistema.
\end{itemize}